% ===========================================================================
% Poste 3 — Assurance-emploi
% ===========================================================================
\poste{3}{Assurance-emploi (AE)}{CA\_ae}

\pdesc Cotisation fédérale sur le revenu d'emploi finançant les prestations
en cas de perte d'emploi. Le Québec bénéficie d'un \emph{taux réduit}, car
son régime parental (RQAP, poste~2) remplace les prestations de maternité et
parentales fédérales.

\pobj Offrir un revenu temporaire aux personnes qui perdent leur emploi.

\pregle Même structure que le RQAP --- pas d'exemption, un seuil et un
plafond :
\[
\text{cotisation} =
\begin{cases}
0 & \text{si } r \le s \\
\tau \times \min(r,\, MRA) & \text{sinon}
\end{cases}
\]
où $MRA$ est le maximum de la rémunération assurable et $\tau$ le taux
réduit du Québec (part de l'employé). Le seuil $s$ reproduit l'article
96(4) de la \emph{Loi sur l'assurance-emploi} : la personne dont la
rémunération assurable de l'année est de \D{2000} ou moins obtient le
remboursement de ses cotisations --- le modèle les met directement à zéro.
La cotisation du ménage est la somme des cotisations des adultes ; les
retraités ne cotisent pas. Au fédéral, la cotisation AE donne droit à un
crédit non remboursable (poste~19).

\emph{Point de contrôle.} Cotisation maximale 2025 :
$0{,}0131 \times \D{65700} = \D{860.67}$ ; 2026 :
$0{,}013 \times \D{68900} = \D{895.70}$.

\pparams

\begin{center}
\small
\begin{tabular}{l r r}
\toprule
Paramètre & 2025 & 2026 \\
\midrule
Maximum de la rémunération assurable ($MRA$) & \D{65700} & \D{68900} \\
Seuil de remboursement ($s$)                 & \D{2000}  & \D{2000} \\
Taux de cotisation, Québec ($\tau$)          & \pc{1.31} & \pc{1.30} \\
\bottomrule
\end{tabular}
\end{center}

% ============================================================================
% GÉNÉRÉ par scripts/exemples-guide.ts — NE PAS ÉDITER À LA MAIN.
% Régénérer : npm run exemples (chaque chiffre est vérifié contre le moteur).
% ============================================================================
\pexemples Même structure que le RQAP : taux unique (taux réduit du Québec,
\pc{1.31}), sans exemption, plafonné au maximum de la rémunération
assurable (\D{65700}). Sous \D{2000} de rémunération, la cotisation est
remboursée (art.~96(4) LAE).

\medskip\noindent\textbf{M2 --- personne vivant seule, 30~ans, revenu de travail de \D{9000}.}
\begin{align*}
\text{cotisation} &= \num{9000} \times \pc{1.31} = \num{117.90}
\end{align*}
Cotisation AE : \D{117.90}.

\medskip\noindent\textbf{M4 --- personne vivant seule, 40~ans, revenu de travail de \D{50000}.}
\begin{align*}
\text{cotisation} &= \num{50000} \times \pc{1.31} = \num{655}
\end{align*}
Cotisation AE : \D{655}.

\medskip\noindent\textbf{M5 --- personne vivant seule, 45~ans, revenu de travail de \D{100000}.}
\begin{align*}
\text{cotisation} &= \min(\num{100000},\ \num{65700}) \times \pc{1.31} = \num{860.67}
\end{align*}
Le revenu excède le MRA : cotisation maximale 2025, \D{860.67}.


\prefs Loi sur l'assurance-emploi (L.C. 1996, ch.~23), art.~4 (maximum),
66 (taux), 69(2) (réduction pour régime provincial) et 96(4)
(remboursement) ; Service Canada ; Bureau de l'actuaire en chef (BSIF).
